\theory{Semigroup theory}{What algebraic structure remains when a composition operation is associative but elements need not have inverses?}{A semigroup is a set with an associative internal operation. It captures processes that can be performed in sequence, even when a process cannot be undone. Semigroup theory studies ideals, Green's relations, representations, finite structure, and links with automata and operators.}{String processing, finite automata, Markov processes, partial differential equations, linear systems, and the algebra of irreversible computation.}{Associative multiplication without universal inverses models irreversible composition, transition systems, and transformation actions.}

\paragraph{Alexander Sushkevich and David Rees}
Sushkevich studied the structure of finite semigroups, while Rees described important completely simple semigroups through matrix-like constructions. Their work showed how associative composition can be analyzed without assuming inverses.

\paragraph{James Alexander Green and automata theorists}
Green introduced the relations that organize semigroup elements by the principal ideals they generate. Later work connected these relations with finite automata, transformation systems, and operator theory, where irreversible composition is the central example \cite{semigroup_ref1}.
\end{historylist}
\section*{3. Theory}
\subsection*{Core theory}
\paragraph{Definition and examples}
A semigroup is a set $S$ with a product satisfying
\[
(ab)c=a(bc)
\]
for all $a,b,c\in S$. No identity and no inverse are required. String concatenation is associative: grouping the joins differently does not change the final string. Functions under composition form a semigroup, as do square matrices under multiplication and positive integers under addition.

A monoid is a semigroup with an identity. A group is a monoid in which every element has an inverse. Thus semigroups generalize groups by allowing irreversible processes such as a reset button, a lossy computation, or a function that is not reversible \cite{semigroup_ref1}.

\paragraph{History and structure}
The systematic study began in the early twentieth century. Anton Sushkevich studied finite simple semigroups; David Rees and James Green developed structural tools; Alfred Clifford and Gordon Preston published foundational monographs. Green's relations divide a semigroup according to the principal ideals generated by its elements \cite{semigroup_ref2}.

An ideal is a subset absorbing multiplication from the relevant side. Minimal ideals describe the long-term core of a finite semigroup. In a finite semigroup, repeated products eventually repeat, so powers of an element enter a periodic part. This finite behavior is useful when a process is applied repeatedly.

\paragraph{Representations and finite semigroups}
Cayley's theorem for semigroups represents every semigroup as a semigroup of transformations of a set. An element is therefore a function, and multiplication is composition. This generalizes the group representation by permutations, replacing reversible bijections with arbitrary maps.

Finite semigroups are closely related to finite automata. A word read by an automaton induces a transformation of its state set. Words with the same effect form a syntactic monoid, and algebraic properties of this monoid characterize language classes. Krohn--Rhodes theory decomposes finite semigroups into simpler group and aperiodic components, analogous in spirit to factoring \cite{semigroup_ref3}.

\paragraph{Analysis and differential equations}
A one-parameter operator semigroup is a family $T(t)$ with $T(0)=I$ and $T(t+s)=T(t)T(s)$ for $t,s\geq0$. It describes evolution forward in time. Unlike a group, negative time may not be possible. If $A$ is the generator, formally $T(t)=e^{tA}$.

The heat equation is an example. Starting with a temperature distribution $u(0)$, the heat semigroup sends it to $u(t)$ for later times. Diffusion smooths irregularities and loses information, so reversing it is unstable. Semigroup methods make existence, uniqueness, and stability of evolution equations precise \cite{semigroup_ref4}.

\section*{4. Important theorems and results}
\begin{enumerate}[leftmargin=*,itemsep=0.35em]
\item \textbf{Cayley representation theorem.} Every semigroup is isomorphic to a semigroup of transformations under composition.

\item \textbf{Green's relations.} The relations $\mathcal L,\mathcal R,\mathcal H,\mathcal D,\mathcal J$ organize elements by their principal left, right, and two-sided ideals.

\item \textbf{Krohn--Rhodes theorem.} Finite semigroups admit decompositions into group components and aperiodic components in a precise wreath-product sense.

\item \textbf{Syntactic monoid theorem.} A regular language has a finite syntactic monoid, and its algebra records the language's behavior under contexts.

\item \textbf{Hille--Yosida principle.} Suitable conditions on an operator characterize when it generates a strongly continuous semigroup.
\end{enumerate}
\noindent\emph{Source for these results:} \cite{semigroup_ref1}.

\theoryapplications
\theoryconnections{semigroup_theory}{%
\item A semigroup keeps the associative composition law while allowing inverses to be absent, which is exactly what irreversible processes require.
\item Functions under composition form the basic example: applying one transition after another gives a new transition, but it may not be reversible.
\item Automata use finite transformation semigroups to record the effect of reading words, while probability and functional analysis use semigroups to describe repeated or continuous evolution \cite{semigroup_ref1,semigroup_ref2}.
}{%
\item In formal-language theory, the syntactic semigroup compresses all words that have the same effect on acceptance, turning an infinite language question into a finite algebraic one when possible.
\item In a Markov process, a transition semigroup records the distribution after each elapsed time.
\item For a PDE such as heat flow, an operator semigroup maps the initial temperature to its later temperature, so composition expresses the passage of time \cite{semigroup_ref1,semigroup_ref3}.
}
\section*{7. Sample Questions}
\emph{Exercise reference:} \cite{semigroup_ref1}. The cited edition is the source for the exercise set; consult its Exercises section for the official statement and numbering.
\begin{enumerate}[leftmargin=*,itemsep=0.9em]
\item \textbf{Exercise 1. Why is string concatenation a semigroup?} Concatenating three strings in either order of parentheses produces the same sequence of symbols, so associativity holds.

\item \textbf{Exercise 2. Are all functions from a set to itself a monoid?} Yes. Composition is associative and the identity function is the identity element. It is generally not a group because most functions are not invertible.

\item \textbf{Exercise 3. What is the identity in the matrix semigroup of all $2$ by $2$ matrices?} The identity matrix is an identity for multiplication, but the set is not a group because singular matrices have no inverse.

\item \textbf{Exercise 4. Why is a heat evolution naturally a semigroup rather than a group?} Heat moves forward and smooths information. A later temperature profile usually does not determine a stable earlier profile, so backward evolution is not available as a bounded inverse.

\item \textbf{Exercise 5. How does an automaton produce a semigroup?} Each input word acts on the finite state set by a transformation. Combining words concatenates their transformations, giving a semigroup under composition.

\end{enumerate}
\section*{8. References}
\begin{thebibliography}{9}
\bibitem{semigroup_ref1} A. H. Clifford and G. B. Preston, \emph{The Algebraic Theory of Semigroups}.
\bibitem{semigroup_ref2} J. M. Howie, \emph{Fundamentals of Semigroup Theory}.
\bibitem{semigroup_ref3} J.-E. Pin, \emph{Mathematical Foundations of Automata Theory}.
\bibitem{semigroup_ref4} E. B. Davies, \emph{One-Parameter Semigroups}.
\end{thebibliography}
